\documentclass[10pt,a4paper]{article}
\usepackage[margin=24mm]{geometry}
\usepackage{amsmath,amssymb,amsthm}
\usepackage[T1]{fontenc}
\usepackage{lmodern}
\usepackage[hidelinks]{hyperref}
\setlength{\emergencystretch}{3em}
\newtheorem{theorem}{Theorem}
\newtheorem{lemma}{Lemma}
\title{A Trivial Picard--Vessiot Extension\protect\\Contradicts Conjecture 00000008502}
\author{idealistichacker\protect\\\small AI-assisted proof; no affiliation claimed}
\date{14 October 2026}
\begin{document}
\maketitle
\begin{abstract}
The conjecture asserts that the Zariski dimensions of the differential Galois
groups of equations $y''=ay$ form exactly $\{1,2,3,4\}$. Over
$K=\mathbb C(t)$ with its usual derivation, $a=0$ gives the Picard--Vessiot
extension $L=K$. The fundamental matrix has determinant one and entries already
in the base field. Thus the actual differential Galois group is trivial;
its represented Zariski closure is the identity point and has coordinate ring
isomorphic to $\mathbb C$, hence dimension zero. This contradicts the asserted
exact spectrum. The Lean development constructs rational-function differentiation,
its constant field, the fundamental matrix, the self-extension automorphism group,
the representation, Zariski closure and coordinate-ring dimension.
\end{abstract}

\section{Source and scope}
Both source languages assert an exact dimension spectrum $\{1,2,3,4\}$ for
second-order equations $y''=ay$, with every listed dimension realized. This is
one necessary clause of a conjunction also claiming a general $n^2$ upper
bound and an equality classification. Refuting the exact-spectrum clause
suffices; this paper does not separately refute the general upper bound.

We use ordinary characteristic-zero Picard--Vessiot theory with algebraically
closed constants. The source does not exclude $a=0$, reducible equations,
trivial groups or identity extensions. This counterexample satisfies the
standard restrictions just stated. It does not require a speculative choice
of a parameterized or analytic-monodromy Galois group. The standard fundamental
matrix and PV-field definitions are described in Singer~\cite{Singer},
Definitions 1.2.13 and 1.3.4 and the PV-field discussion preceding the latter.

The fixed original is preserved separately; it is not edited.
\begin{description}
\item[Source commit] \hfill\\
\texttt{cce2d90cb298de9595798eb674d2c80712100532}
\item[Conjecture blob] \hfill\\
\texttt{a2f2c93606d87db922691efc471572a32355e881}
\end{description}

\section{The concrete differential field}
Take $K=\mathbb C(t)$ and differentiate rational functions by
\[
D\!\left(\frac{p}{q}\right)=\frac{p'q-pq'}{q^2}\qquad(q\ne0).
\]
The usual cancellation calculation shows this is independent of the chosen
fraction representation, additive, and satisfies the Leibniz rule. It kills
every coefficient from $\mathbb C$, and $D(t)=1$.

\begin{lemma}
The constants of $K$ are exactly $\mathbb C$.
\end{lemma}
\begin{proof}
For nonzero reduced $p/q$, if its derivative vanishes then $p'q=pq'$.
Because $p$ and $q$ are coprime, $p$ divides $p'$. A nonzero polynomial
derivative has strictly smaller degree, so $p'=0$. Now $p\ne0$ implies
$q'=0$. In characteristic zero, both $p$ and $q$ are constant polynomials.
The zero rational function is also a constant. Conversely differentiation
kills all constant coefficients.
\end{proof}

These properties are proved on Mathlib's actual \texttt{RatFunc} type,
including its genuine coprime numerator and denominator. The derivative is
not a label attached to an arbitrary function. Algebraic closedness of the
coefficient field uses Mathlib's fundamental theorem of algebra.

\section{A Picard--Vessiot self-extension}
For $a=0$ the companion system and fundamental matrix are
\[
A=\begin{pmatrix}0&1\\0&0\end{pmatrix},\qquad
U=\begin{pmatrix}1&t\\0&1\end{pmatrix}.
\]
Entrywise differentiation gives $D(U)=AU$, and $\det U=1$. Its columns are
the scalar solutions $1,t$ with their first derivatives. The nonzero
Wronskian makes this an invertible fundamental matrix for the order-two
equation, not a scalar equation of reduced order.

All entries of $U$ and $1/\det U$ already lie in $K$. The ring generated
over $K$ by these entries is exactly $K$, which is differentially simple
because a field has only the ideals $0$ and $K$. Its fraction field is $K$.
Consequently $L=K$ is the generated PV field and is minimal. The extension
has no new constants, since it is the identity extension; its constants
are exactly the $\mathbb C$ verified above. Properness of a PV extension
is not part of the standard definition.

\section{The actual Galois group and algebraic dimension}
Let $G=\operatorname{Aut}_D(L/K)$, the derivative-commuting subgroup of
the actual $K$-algebra automorphisms of $L$. At $L=K$, an automorphism fixing
the base fixes every element. Thus $G$ has exactly the identity automorphism.
The standard representation determined by the fundamental matrix is
\[
\rho(\sigma)=U^{-1}\sigma(U)=I_2.
\]
Its entries belong to the constant field. Lean checks the displayed formula
against the constructed homomorphism, not just against a declared trivial-group
label.

Represent $\mathrm{GL}_2(\mathbb C)$ by five affine coordinates
$(x_{00},x_{01},x_{10},x_{11},z)$ satisfying
$z(x_{00}x_{11}-x_{01}x_{10})=1$, so $z$ is the inverse determinant.
The identity point is $(1,0,0,1,1)$. It is Zariski closed: coordinate-equality
polynomials cut out precisely this point. Hence the represented group's
Zariski closure is still this point.

Its coordinate ring is the actual polynomial ring quotient by the point's
vanishing ideal. Evaluation at the identity is surjective, and its kernel
is exactly that vanishing ideal. The first isomorphism theorem identifies
the quotient ring with $\mathbb C$. The Krull dimension is therefore zero.
This argument does not replace algebraic dimension by the cardinality of a
group, or assign a dimension by definition.

\begin{theorem}
The exact second-order spectrum asserted in conjecture 00000008502 is false.
\end{theorem}
\begin{proof}
The verified equation $y''=0$, base $K=\mathbb C(t)$ and PV field $L=K$
give an admissible actual represented group of Zariski dimension zero.
But $0\notin\{1,2,3,4\}$. Thus this exact-spectrum necessary clause, and
therefore the source conjunction, cannot hold.
\end{proof}

\section{Lean correspondence and reproducibility}
The source correspondence is:
\begin{description}
\item[\texttt{RationalDerivative.lean}] \hfill\\
Constructs the derivative, its laws and the exact constants.
\item[\texttt{DifferentialMatrix.lean}] \hfill\\
Proves the matrix and self-automorphism components.
\item[\texttt{PointClosure.lean}] \hfill\\
Implements polynomial zero loci, the induced standard Zariski topology,
closure, invertible-matrix coordinates and the genuine coordinate quotient.
\item[\texttt{ConcretePV.lean}] \hfill\\
Assembles the concrete PV data, actual represented group and closed negation
of the universal exact-spectrum necessary clause.
\end{description}

All 76 public theorem types and transitive axioms are printed by
\texttt{ResearchCheck.lean}. Only the standard axioms \texttt{propext},
\texttt{Classical.choice} and \texttt{Quot.sound} are allowed. The project uses
Lean 4.33.1 and Mathlib commit
\texttt{0df444a360eaa60ab8c11dca51a86af692955474}. The installed compiler and
fixed library remain trusted infrastructure; no full compiler-source audit is claimed.

The research reproducer builds all authored modules in a fresh private directory,
audits every listed theorem and runs exact rational regression checks. It reuses
dependency caches: it is not a full dependency rebuild or empty-machine bootstrap.
Finite regressions support reproducibility, not the universal mathematical argument.

This document is preparation material, not a claim of submission, first-solver
status, human endorsement, upstream acceptance, ranking or award. A portable
submission package, final frozen independent review, actual PDF reproduction and
fresh publication gates remain separate requirements.

\begin{thebibliography}{1}
\bibitem{Singer}
Michael F. Singer, \emph{Introduction to the Galois Theory of Linear
Differential Equations}, arXiv:0712.4124v2, 2008.
\end{thebibliography}
\end{document}
